% !TeX root = ../papers/nc-p3.tex
\section{四つの同方向の三行解と長さ十七の固定点スキーム}
\label{sec:directed-binary-defect-cycle}
\begin{theorem}\label{thm:directed-binary-defect-cycle}
系\ref{cor:two-dimensional-full-arrow-apolarity}の階数4の四矢作用を持つ二次元印付き有限窓を考える。
各位置で命題\ref{prop:two-dimensional-quadric-successors}の後続行空間が三次元とする。
この印付き層には次のパラメータ表示がある。四本の次数2の射
$g_i:P_i\to P_{i+1}$（$P_i=\mathbb P(U_i)$）を取り、$D$ を
$x_{i+1}=g_i(x_i)$ の巡回交叉とする。任意の一座標への射影は $D$ を
四射の合成の固定スキーム（長さ17）に同型に写す。
$H^0(D,\mathcal O_D(2,2,2,2))^*$ の射影点を一つ選ぶ。その汎関数の最小支持部分スキームの
長さは $15,16,17$ のいずれかに限れる。

各 $i$ で $x_i,x_{i+2}$ 上の九次二元形式の十次元空間同士を
$D$ 上で掛けて汎関数を取った十行十列の対合を作る。ここで第一因子の線束は
$(1,2,1,0)$ の巡回移動、反対因子はその補次数の制限である。
反対二位置は転置なので独立に検査するのは二行列で足りる。
両階数が6で、復元テンソルが標準有限Frobenius条件と四辺の有限性を満たす点を採用する。
これでこの候補層を全て、非被約な固定点を含めて得る。
同型判定は四つの射と汎関数の直線を同時に運ぶ $\prod_i\mathrm{PGL}(U_i)$ の作用である。
逆向きの列空間四つについても添字の順を逆にした表示が成り立つ。
\end{theorem}
\begin{proof}
文字を $\operatorname{Hom}(U_i,U_{i+1})$ とした全作用表示で、
位置 $i-1$ からの後続行は $V_i^*\otimes U_i^*=\operatorname{End}(U_i)\otimes U_{i+1}^*$ に属する。
既知の二行はスカラー行列と $U_{i+1}^*$ のテンソル積である。
体積形式を選ぶとこの二行を除いた空間は $\operatorname{Sym}^2 U_i^*\otimes U_{i+1}^*$ になる。
行の方程式はこの $(2,1)$ 形式による $F$ の縮約が零になることに等しい。
実際文字 $i$ の出口で一因子、入口で前文字の出口も合わせ二因子を縮約する。
交代成分は $F$ の対称性で消え、これがちょうど先のスカラー二行である。
余分な行は倍率まで一意なのでその $(2,1)$ 形式を $f_i$ とする。

$U_{i+1}^*$ の基底について余分な行を跡零行列の組 $B_0,B_1$ と書く。
対応する二次式は $\det(u,B_j u)$ である。両式が同じ $[u]$ で消えれば、
各 $B_j$ からその固有値を引いて両方が $u$ を殺すようにできる。
この行は既知の二行を足した非零の行解で、$U_i$ 上の行階数が1になり、有限辺に反する。
従って二式に共通零点はなく、$f_i=0$ は次数2の射のグラフである。
四グラフの交叉は逐次代入により、任意の座標上で次数16の射の固定点に等しい。
その射は恒等になれないから長さは17であり、四方程式は滑らかな四次元積で正則列を局所的に与える。

$L=\mathcal O(2,2,2,2)$ とする。四グラフのKoszul分解を $L$ で捻ると、
選んだ方程式を示す $e_i\in\{0,1\}$ に対して各位置の次数は $2-2e_i-e_{i-1}\ge-1$。
よって全項に高次コホモロジーがなく、$H^0(L)\to H^0(D,L)$ は全射で、
核は $f_i$ と補次数の形式の積で張られる。縮約条件は正確にこの核の零化である。
従って一つの汎関数から一つの $F$ が戻り、逆も一意である。逆構成で有限辺と指定階数を課せば
各非零 $f_i$ は真に第三行を与え、四つの射の一意性も従う。全作用表示の基底変更は前掲の系の通りである。

左切断は $P_i$ 上で $1+2\cdot2+4=9$ 次となり、得られる線型系は九次の全体である。
どの次数2の射 $g$ でも $H^0(\mathcal O(1))\,g^*H^0(\mathcal O(m))=H^0(\mathcal O(2m+1))$ である。
実際基点のない二次式の鉛筆は標数零では左右の射影変換で $[s^2:t^2]$ になる
（二重根の二つの二次式を選ぶ。判別二次曲線への接線なら共通因子を持つ）。
この座標で偶奇の単項式を取ればよい。まず次の射の二次式鉛筆に全二次式を掛けて全四次式を得る
（同じ座標で指数 $0,1,2$ と $0,2$ を足す）。それを一つ手前に引き戻して一次式を掛ければ主張を得る。
右切断も反対座標上の全九次式になる。

汎関数の局所積対合の零化イデアルを割った部分スキームを $D'$、その長さを $r$ とする。
各局所環は一変数の打切り環なので、$D'$ は最小支持で、そこでの積対合は完全である。
両九次式の評価階数は $\min(10,r)$ であり、合成の階数は少なくとも $2\min(10,r)-r$。
$r\le10$ の場合は階数が $r$ そのものだから $r=6$ だが、既出の長さ6の純関係障害で不可能である。
従って $r\ge14$。等号をさらに排除するため、$x=x_i,y=x_{i+2}$ とし、
$D'$ 上で $y=p(x)$、$p$ は次数4の射であることを使う。
$r=14$ では九次式の対合の左核は次元4であり、その $D'$ 上の評価は
右九次式の完全対合における零化空間（次元4）に一致する。後者は $y$ 上の全三次式を
一つの可逆な局所倍率で掛けた評価である。
実際無限遠を避けて $D'$ をモニック十四次式 $d(y)$ で表示し、$k[y]/d$ の積の十三次係数を取ると
完全対合になり、九次以下の直交は三次以下である。任意の他の完全な積汎関数は単元倍である。

従って左核の四つの九次式 $a_j(x)$ を、$D'$ での射影値が $y$ の三次Veroneseとなるように選べる。
$p=[p_0:p_1]$ とすれば十三次式 $p_0a_{j+1}-p_1a_j$ ($0\le j<3$) は
長さ14の $D'$ で消え、恒等的に零である。互いに素な四次式 $p_0,p_1$ に対し、
これらの等式は $a_0$ が $p_0^3$ で割れることを意味し、$9<12$ に矛盾する。
この議論は固定点の分離性を用いない。
\end{proof}
